\section{作品概述}

\subsection{背景分析}

\subsubsection{先理解智能体：会行动的软件}

普通聊天模型的主要输出是一段文字；智能体则把大模型、外部内容、持久化记忆和工具调用组合成一个能够持续完成任务的执行系统。它可以读取网页或知识库，记住上一次任务的结果，根据模型判断选择工具并填写参数，最后发送邮件、读写文件、调用网络 API 或执行代码。因此，智能体的安全边界不只在模型输入和输出，还在模型把意图转换为真实动作的那一刻。

智能体与传统信息系统有一个关键差别：传统程序通常按照预先编写的路径调用函数，而智能体会根据自然语言和外部内容动态选择工具、组织调用顺序并生成参数。外部内容因此不再只是被动的数据，它可能影响后续计划；模型输出也不再只是给人阅读的文本，它可能直接成为工具请求。对于信息安全人员而言，这意味着需要同时保护\textbf{内容边界、决策边界和执行边界}，而不能只判断一句话是否“看起来有害”。

\begin{figure}[htbp]
\centering
\reportFigureLead{不可信内容会改变计划与工具参数；运行时屏障在具体攻击后果发生前执行最后一次强制判定。}
\resizebox{0.98\linewidth}{!}{%
\begin{tikzpicture}[
  font=\small,
  box/.style={draw, rounded corners=2pt, minimum width=2.75cm,
              minimum height=1.35cm, align=center, fill=reportSystem!10, draw=reportSystem},
  consequence/.style={draw, dashed, rounded corners=2pt, minimum width=3.4cm,
                      minimum height=1.55cm, align=center, fill=reportBlock!10, draw=reportBlock},
  barrier/.style={draw, very thick, rounded corners=2pt, minimum width=3.55cm,
                  minimum height=1.55cm, align=center, fill=reportEvidence!13, draw=reportEvidence},
  a/.style={-{Stealth[length=1.8mm]}, thick},
  bad/.style={-{Stealth[length=1.6mm]}, dashed, thick, draw=reportBlock},
  stop/.style={-{Stealth[length=1.8mm]}, thick, draw=reportPass}
]
\node[box] (input) at (0,0) {\textbf{1 不可信内容}\\检索、工具返回\\持久化记忆};
\node[box] (agent) at (3.45,0) {\textbf{2 模型规划}\\选择工具\\组织调用顺序};
\node[box] (params) at (6.9,0) {\textbf{3 工具参数}\\收件人、文件路径\\API 与消息载荷};
\node[consequence] (harm) at (10.75,0) {\textbf{4 具体攻击后果}\\邮件外送敏感数据\\越权读文件 / 记忆投毒};
\node[barrier] (guard) at (15.15,0) {\textbf{5 最后责任点}\\运行时屏障先于副作用\\deny / ask / allow\_sanitized};

\draw[a] (input) -- (agent);
\draw[a] (agent) -- (params);
\draw[bad] (params) -- (harm) node[midway,above,font=\scriptsize] {若直接执行};
\draw[stop] (harm) -- (guard) node[midway,above,font=\scriptsize] {后果发生前拦截};
\node[below=0.35cm of guard, font=\scriptsize\bfseries, text=reportPass]
  {阻断或脱敏后才进入真实工具；原始危险调用执行 0 次};
\end{tikzpicture}%
}
\caption{不可信内容到运行时拦截的五步因果链。模型规划会形成具体工具参数；若直接执行将产生数据外送、越权读取或记忆投毒，本作品在最后责任点先完成强制判定。}
\label{fig:problem-role}
\end{figure}

\begin{center}
\reportOptionalImage{figures/gpt-agent-action-risk.png}{展示“不可信外部内容—模型规划—工具参数—真实后果—运行时拦截”的介绍性场景；提示词见 \texttt{visual-prompts.md}。}
\end{center}

\subsubsection{问题是什么：不可信内容会被转成真实动作}

攻击者不必让模型直接生成一段明显有害的文字，只需把恶意意图嵌入检索结果、工具返回值或持久化记忆，就可能改变智能体的下一步计划。典型链路是：不可信内容进入上下文，模型据此提出工具调用，工具适配层把调用转换成参数，最后由外部系统执行。只要其中一个环节把“模型提出的动作”误当成“可以直接执行的动作”，攻击就会产生真实副作用。

\begin{table}[htbp]
\centering
\caption{智能体运行时安全问题、具体后果与所需控制}
\label{tab:agent-runtime-problems}
\small
\begin{tabular}{@{}p{2.55cm}p{3.6cm}p{3.45cm}p{3.25cm}@{}}
\toprule
问题来源 & 攻击者如何影响智能体 & 可能造成的直接后果 & 需要的安全控制 \\
\midrule
检索内容或工具返回值 & 在看似正常的资料中加入改变工具参数的指令 & 外送内部邮件、调用错误 API、把敏感内容交给第三方 & 在工具执行前重新检查来源、意图和参数 \\
模型输出与调用计划 & 让模型生成绕过策略的工具请求或危险步骤 & 读取受保护文件、执行越权操作、触发代码副作用 & 在模型输出交付前和工具请求前设置独立屏障 \\
持久化记忆或任务状态 & 写入会影响后续任务的恶意事实或指令 & 后续会话持续执行错误策略，形成记忆投毒和权限扩大 & 对写入和读取同时进行运行时监督 \\
外部裁判或远程服务 & 把原始内容、来源和真实参数一并发送出本地边界 & 私有数据泄露，且泄露后无法由本地工具撤回 & 先确定性脱敏，只外送最小必要载荷 \\
\bottomrule
\end{tabular}
\end{table}

本作品用“不可信检索内容劫持 \verb|send_email|”作为可复现的代表攻击链：检索结果没有出现明显危险词，模型仍然提出邮件调用；如果系统只在文本层检查，邮件参数可能已经被改写；如果系统在工具执行点重新核验，攻击链就会在邮件外送前停止。这个例子说明，作品要解决的不是“让模型永远不犯错”，而是\textbf{即使模型被外部内容影响，也不能把未经过安全判定的动作交给真实工具执行}。

\subsubsection{问题为什么重要：一次错误动作可能不可逆}

智能体运行时安全的重要性来自三个相互叠加的事实。第一，攻击内容可以伪装成正常的业务资料、网页片段或历史记忆，使用者很难仅凭阅读模型对话发现它已经改变了执行计划。第二，工具动作往往具有不可逆的副作用：邮件一旦发出无法收回，敏感文件一旦被读取可能已经泄露，写入长期记忆后还会影响后续任务。第三，智能体能够自动重复调用工具，一个很小的判断偏差可能在无人值守的任务链中被放大为多次外送、批量读写或持续性配置修改。

因此，这不是单纯的模型质量问题，而是一个典型的信息安全控制问题：需要确认\textbf{谁发起了动作、模型判定的对象是否就是即将执行的对象、动作是否越过了数据边界，以及系统故障时是否会默认放行}。这些问题分别对应身份绑定、参数完整性、隐私边界和失效闭合，必须在运行时而不是事后日志中解决。

\subsubsection{现有防护为什么仍需要运行时安全屏障}

模型对齐和提示词过滤可以降低明显的越狱或有害请求，但它们通常观察的是输入和文本输出；它们不能单独保证模型最终生成的工具名称、调用标识和参数没有在上下文中被悄悄改变。静态工具权限可以限制“能不能调用某个工具”，却不一定能判断“这一次调用的对象、参数和来源是否符合当前任务”。容器、虚拟机等资源隔离可以降低文件系统或进程层面的影响，但无法替代邮件外送、远程 API 和消息交付前的语义判定。

本作品因此把安全控制放在真实动作边界：运行前检查上下文和任务，模型输出交付前检查生成结果，工具执行前检查真实参数，消息交付前检查最终外送内容。OWASP、MITRE ATLAS、CSA AI Controls Matrix 与 NIST AI RMF 对提示注入、工具滥用、配置暴露、可追溯性和失效处理等风险的归纳，为这一运行时控制面提供了行业参照\cite{owaspllm2026,mitreatlas2026,csaaicm2026,nistairmf2023,nistgenai2024}。与本作品直接相关的 MITRE ATLAS 技术如下。

\begin{table}[htbp]
\centering
\caption{MITRE ATLAS 中与智能体运行时安全直接相关的技术（版本 2026.07）}
\label{tab:atlas}
\small
\begin{tabular}{@{}>{\raggedright\arraybackslash}p{1.95cm}>{\raggedright\arraybackslash}p{4.85cm}>{\raggedright\arraybackslash}p{5.5cm}@{}}
\toprule
技术编号 & 技术名称 & 战术归属 \\
\midrule
AML.T0084 & Discover AI Agent Configuration & 发现智能体配置 \\
\quad .000 & \quad Embedded Knowledge & \quad 发现可访问的数据源 \\
\quad .001 & \quad Tool Definitions & \quad 发现可用工具 \\
\quad .002 & \quad Activation Triggers & \quad 发现激活触发条件 \\
\quad .003 & \quad Call Chains & \quad 提取调用链，定位远程代码执行目标 \\
AML.T0081 & Modify AI Agent Configuration & 篡改配置以持久化 \\
AML.T0083 & Credentials from AI Agent Configuration & 从配置中获取其他工具凭据 \\
AML.T0097 & Virtualization/Sandbox Evasion & Defense Evasion（AML.TA0007） \\
AML.T0105 & Escape to Host & Privilege Escalation（AML.TA0012） \\
\bottomrule
\end{tabular}
\end{table}

\subsubsection{本作品解决什么问题：在动作发生前建立可执行的安全屏障}

\begin{center}
\setlength{\fboxsep}{8pt}
\fbox{\parbox{0.91\linewidth}{\textbf{一句话概括：}本作品在不替换原有模型和工具的前提下，把不可信内容到真实动作之间缺失的安全决策层补上；它先判断风险，再决定动作能否继续，并在必要时阻断、询问、脱敏和留证。}}
\end{center}

作品交付的是\textbf{可复用于不同智能体框架的运行时安全中间件}。核心引擎以统一请求、决策、检测器台账和内容无关审计为接口，框架适配层只负责把本框架的生命周期映射到真实执行点。OpenClaw 已完成首个接入；其他框架不需要重写安全策略，只需实现相同的生命周期映射。“沙箱”在此指对来源、意图和动作实施强制策略的语义与执行屏障，可以与容器、虚拟机等资源隔离层协同，而不是用一个容器替代全部安全控制。

本作品把问题拆成四个必须同时成立的安全性质：

\textbf{第一，明确高危必须立即停止。}规则层命中高危信号后直接短路，后续模型和外部裁判不再获得调用机会；这既减少攻击面，也避免把已经确定的危险内容继续交给更大的处理链。

\textbf{第二，模糊风险可以升级但不能扩大泄露。}只有前一级无法解决的信号才进入下一级；需要外部裁判时，先在本地确定性脱敏和令牌化，原始内容、来源句柄、真实工具名和原始参数不得直接外送。

\textbf{第三，被评估的对象必须就是将要执行的对象。}OpenClaw 工具点同时绑定用户输入、模型输出投影和真实工具请求，对调用标识、工具名和参数做深比较，阻断“评估一个调用、实际执行另一个调用”。

\textbf{第四，任何结果都必须转化为可执行动作。}允许、询问、拒绝、脱敏后允许四值决策直接控制执行屏障；询问和拒绝在副作用发生前冻结工具，脱敏后允许只把变换后的载荷交给工具，检测器失败、预算耗尽或执行层健康失败均失效闭合为拒绝或询问，而不是默认放行。

\subsection{相关概念}

参考报告通常先定义作品中会反复出现的概念，再进入详细设计。本作品面向的信息安全读者不一定从事大模型或智能体研究，因此本节只解释阅读后文所必需的术语；具体实现、参数和代码仍放在第二章和第三章。

\subsubsection{大模型智能体}

\textbf{大模型}是能够根据上下文生成文字、结构化结果或下一步决策的模型；\textbf{智能体}是在大模型外增加了任务循环、外部数据、持久化记忆和工具调用的软件系统。模型负责理解和规划，工具负责把规划变成读写文件、发送消息、访问网络或执行代码等动作。本报告中的“智能体”特指这种能够继续执行动作的组合系统，而不是只返回文字的聊天窗口。

\subsubsection{提示注入、检索投毒与记忆投毒}

\textbf{提示注入}是攻击者把改变模型行为的指令放进用户提示或其他输入中；\textbf{检索投毒}是把类似指令藏进网页、知识库或工具返回值，使模型把不可信资料当成任务要求；\textbf{记忆投毒}是把恶意事实或规则写进持久化记忆，使后续任务继续受到影响。三者的入口不同，但共同结果都是让模型的计划或工具参数偏离原任务，因此需要在真实动作前统一监督。

\subsubsection{工具调用与真实执行点}

\textbf{工具调用}是模型提出的一条结构化请求，通常包含调用标识、工具名称和参数；它本身还不等于动作已经发生。\textbf{真实执行点}是工具、文件、网络或消息真正产生副作用前的 awaited 屏障。本作品在 OpenClaw 中设置四个真实执行点：运行前、模型输出交付前、工具执行前和消息交付前。只有通过屏障的调用才会继续进入原有业务流程。

\subsubsection{运行时安全中间件与语义沙箱}

\textbf{运行时安全中间件}位于智能体框架和模型、工具之间，接收统一的评估请求，返回统一的安全决策，并把决策映射回原框架的执行流程。它不替换模型，也不把工具实现重新写一遍，而是在动作发生前增加一层可复用的检查和控制。本文所说的\textbf{语义沙箱}不是容器或虚拟机本身，而是对内容来源、模型意图、工具参数和外送动作实施策略的屏障；它可以与文件系统、进程和网络层的资源隔离协同工作。

\subsubsection{三级级联与四值动作}

\textbf{三级级联}是按风险确定性逐级增加判断能力的处理链：第一级是确定性规则，第二级是本地模型，第三级是经过脱敏的外部裁判。高危信号在第一级立即短路，只有未解决的信号才升级，避免把明确危险内容继续交给更多处理环节。\textbf{四值动作}是统一的终局结果：允许（allow）继续原调用，询问（ask）冻结副作用并等待操作员裁量，拒绝（deny）直接阻断动作，脱敏后允许（allow\_sanitized）只把确定性变换后的载荷交给工具。

\subsubsection{确定性脱敏与令牌化}

\textbf{脱敏}是从原始内容中移除或替换不应离开本地边界的字段；\textbf{令牌化}是用不包含原文的标识代替来源句柄、工具参数等敏感值。本文的“确定性”表示：在版本和配置相同的情况下，同一输入得到相同的字段集合和替换结果。外部裁判只能看到完成边界校验的最小必要载荷，不能直接看到原始正文、真实工具名和原始参数。

\subsubsection{内容无关审计与 P6/P7 证据}

\textbf{内容无关审计}只记录风险类别、决策动作、执行点状态和哈希引用，不把用户原文复制进日志；这样既能复核判定过程，也不会让审计系统成为新的泄露渠道。P6 是评测结构完整性证据，关注捕获、封印、签名收据和重放信封是否齐全；P7 回放信封中封印的 provider 响应，并重新执行安全引擎与评测器，以复现决策树哈希和指标。二者分别回答“证据结构是否完整”和“产物能否确定性重放”；P7 不验证 LLM 推理确定性。

\subsection{作品定位与核心优势}

本作品不是离线的风险分类器，也不是只对模型输出打标签的审计脚本，而是把\textbf{判定、强制执行和证据复核}放在同一条运行时链路上的安全中间件。它把“如何判”交给三级升级式级联，把“如何阻断”交给 OpenClaw 四个真实执行点，把“如何证明确实发生过这些判定”交给 300 条签名哈希绑定的评测证据。

\textbf{三级级联。}规则层首先执行，高危命中立即短路；仅将未解决信号升级到本地模型，仍未解决时才在确定性脱敏后调用外部裁判，任何故障均失效闭合。四值动作把风险判定直接映射为执行控制，避免“检测结果存在但工具照常执行”。

\textbf{真实执行点。}OpenClaw 的运行前、模型输出交付前、工具执行前和消息交付前均设置 awaited 屏障。判定发生在真实副作用之前，而不是离线生成一条事后结论；邮件参数劫持在工具点被拒绝，工具执行保持为零。

\textbf{300 条可复现证据。}夹具、补丁、输入与结果树均由哈希锁定；P6 记录捕获、封印、签名收据和重放信封的完整结构，P7 回放封印的 provider 响应并重新执行安全引擎与评测器，复现决策树哈希，使产物级结果可以逐条复核。

\subsection{应用前景分析}

\subsubsection{开源智能体应用的安全加固}

OpenClaw 证明了中间件的第一种落地方式：目标应用保留原有模型与工具，只在四个执行点增加适配层和统一判定调用。对检索投毒、工具调用劫持和记忆投毒等攻击，监督层在工具真实执行前给出拒绝或询问；对无风险请求，四值动作保留原有业务流程。

\subsubsection{企业内部智能体的合规审计}

企业场景需要同时解决“能否阻断”和“能否留证”。内容无关审计只保留判定、动作、风险等级、类别计数、执行点结果和哈希引用，不落原始正文。外部裁判所需载荷由本地确定性脱敏器生成，原始来源句柄、真实工具名、原始参数和内部定位信息不外送；若载荷结构、令牌或字段集合不符合契约，系统以 \verb|external_redaction_failed| 失败闭合。

\subsubsection{跨框架运行时安全基础设施}

统一契约使同一套级联和审计逻辑可以服务于不同智能体框架。框架适配层只负责把本地生命周期映射为四类执行点，并把真实工具调用绑定到评估请求；策略、检测器路由、脱敏和证据封印留在中间件内部。300 条夹具和 P6/P7 重放流程还可以直接作为其他框架接入后的回归基线。

\subsubsection{与现有纵深防御的协同}

\enlargethispage{3\baselineskip}

运行时中间件补上了模型对齐和虚拟化隔离之间的语义边界：仅靠训练阶段的模型对齐，不能对运行时后来注入的检索内容施加强制处置；虚拟化隔离也无法判断一次格式合法的邮件请求是否受投毒内容驱动。三级级联则在工具、文件和消息动作真正落地前检查这种因果关系。作品因此可以作为现有基础设施的监督层，而不是替代基础设施隔离。
